% $RUN6 = paper/data/run9-visible-exposure
\section{Limitations and Ethics}
\label{sec:limitations}

% src: main.tex Section 1 (the first specification and the main run); $RUN6/README.md (the two defaults changed on 14 November 2026 after a review of results); tables/fmt-base-density.tex, tables/verdicts.tex (the first specification beside the main run)
\paragraph{Every number is exploratory.} All of the runs worked from one and the same exploratory split of the headline archive, with the method altered from one run to the next once results had been looked at. The main run itself was changed in that way. The first specification gave each forecast its own row's realized traffic and scored the base rate with a normal density; the main run fixes exposure and gives the base rate a Student-t density, and the first specification is reported beside it (Table~\ref{tab:fmt-base-density} and Section~\ref{sec:verdicts}). The same route produced the form taken by the precedent rung, the way the text classifier reads the reference headline next to the variant, the guard and the selection rule, and the comparison for the selection rule was run over the very rows that give it its score. The choices about how results are read came the same way, among them the event taken at planned precision, the readings by origin and with quarters as the units, the random order matched on arms, and the rules governing how the hypotheses are read. So each number in this paper counts as exploratory, and runs agreeing with each other on a single split does not amount to replication. The confirmatory split of the archive is still unopened.

\paragraph{The archive's notice.}
% src: Section 3, The randomization notice (the notice, its dates and the date rule); $RUN6/backtest-upworthy-20261106-2dc461f52ea4.json tables.aggregate[cell=ALL,rung=text_lr] brier_skill_planned 0.0153 (0.0069, 0.0238) against $RUN6/include-randomization-window/backtest-upworthy-20261114-4bfb354b87e3.json tables.aggregate[cell=ALL,rung=text_lr] brier_skill_planned 0.0047 (-0.0018, 0.0116) (tables/fmt-window.tex); Section 6, abstention (why one origin abstains on every row)
The maintainers of the archive flag a window of dates for randomization problems, and each test beginning within it is left out of the main analysis (Section~\ref{sec:setting}). The rule operates on dates because the file we read carries no flag column, so a randomization problem beyond that window escapes it and any test that is sound but sits inside gets dropped. What the rule leaves behind is a gap that forces a single origin to abstain across all rows, which follows from the rule itself and is not a finding (Section~\ref{sec:abstention-results}). The ablation in Table~\ref{tab:fmt-window} is the one that includes the window. Some readings switch sides under it (Section~\ref{sec:results}), among them the text classifier's skill read at planned precision, whose interval then spans zero. \iftmlr An earlier version of the code predates this date rule.\else Version 1.0.0 of this paper's package predates this date rule.\fi

\paragraph{Scope.}
% src: Section 3, the two archives; Table tab:schema and docs/tutorials/02-your-own-archive.md (the package takes any ledger of concluded contrasts with an effect and a standard error); no report measures any design below
Randomized A/B and A/B/n tests are what the evidence covers, mostly a single publisher's tests of headlines and images. The ledger format requires only an effect, its standard error and per-arm counts, so other online randomized experiments also fit the package, model launches and pricing among them, though they remain untested. Untested as well are designs that randomize stretches of time or regions instead of users, and designs where the treatment is not randomized at all, and the dependence in such designs may invalidate the gates. Studies that are observational fall outside the scope here.

\paragraph{The archives and the design.} Headline tests from a single publisher represent a single kind of change inside a single product. The retailer archive falls short of what any gate needs in size, and it has no dates. So the reading of the protocol rests on a single publisher alone, and any drift across time shows up on that one archive alone. Everything here rests on a single archive of each kind, a single main seed and a single design.

\paragraph{The scores and the stand-ins.} A true effect never gets observed, so scoring always runs against the estimate the test itself produced, with floors and ceilings set alongside. Neither archive records a planned exposure, and the main event is defined on the realized standard error, a quantity from after launch. The points below deal with both, and with three further limits of the comparisons.
\begin{itemize}
% src: $RUN6/paper-extras-upworthy-20261106-2dc461f52ea4.json impressions_effect_correlation corr_log_impressions_abs_z 0.087, n_rows 10903; exposure_realized.definition (the main run fixes every row's exposure at its cell's median over the rows visible at the row's origin)
\item The headline archive records no planned exposure, so the main run gives each row the median exposure of its cell over the rows visible at its origin. A real plan for a test could carry information that this stand-in removes, and these archives cannot say whether it does. Realized traffic does carry some: the log of realized impressions has a correlation of 0.087 with the absolute standardized effect, that is, the effect over its standard error with the sign dropped. That is why the run reading realized traffic is kept as a diagnostic only (Table~\ref{tab:fmt-exposure-fixed}).
% src: $RUN6/paper-extras-asos-20261103-6bb5a15c5ff0.json impressions_effect_correlation corr_log_impressions_abs_z 0.2152, corr_log_run_length_abs_z 0.1522
\item For the retailer archive, exposure is not fixed, and the stand-in is the count of users at the close of a run whose length varies. The log of that count correlates 0.215 with the absolute standardized effect. The logged run length, read by no feature at all, comes out at a correlation of 0.152.
% src: $RUN6/backtest-upworthy-20261106-2dc461f52ea4.json tables.aggregate[cell=ALL] event_rate 0.0854, event_rate_planned 0.0764 (tables/event-definition-upworthy.tex)
\item Defining the main significant-positive event uses the realized standard error, and the same holds for the quantities going into the floor of each score. In Table~\ref{tab:event-definition} the event defined on the planned standard error is scored as well, and there the event rate goes down from 0.0854 to 0.0764.
% src: $RUN6/backtest-upworthy-20261106-2dc461f52ea4.json tables.aggregate[cell=ALL] brier_skill_planned rung=P1 0.0127 (0.0096, 0.0155), rung=P2 0.0133 (0.0100, 0.0163), rung=text_lr 0.0153 (0.0069, 0.0238); $RUN6/include-randomization-window/backtest-upworthy-20261114-4bfb354b87e3.json tables.aggregate[cell=ALL] brier_skill_planned rung=P1 0.0103 (0.0075, 0.0130), rung=P2 0.0108 (0.0077, 0.0139), rung=text_lr 0.0047 (-0.0018, 0.0116) (tables/fmt-window.tex); $RUN6/backtest-asos-20261103-6bb5a15c5ff0.json and $RUN6/backtest-synthetic-20261103-bf4fd286e2b2.json brier_skill_planned equal to brier_skill for every rung
\item On that planned event the unadjusted intervals of P1, P2 and the raw text classifier lie above zero in the main run. With the randomization window included, those of P1 and P2 stay above zero and the classifier's spans zero (Section~\ref{sec:ladder-results}; Table~\ref{tab:fmt-window}). On the other two archives the two events coincide.
% src: src/mootcourt/predictors/p0.py _CellSeries (cumulative statistics over every visible row, no period term); Section 4 (P1 takes the level of the last complete quarter); $RUN6/backtest-upworthy-20261106-2dc461f52ea4.json tables.aggregate[cell=ALL,rung=P1] brier_skill_planned 0.0127 (0.0096, 0.0155); $RUN6/ablations/backtest-upworthy-20261106-2dc461f52ea4--base-recency-quarter.json tables.aggregate[cell=ALL,rung=P1] brier_skill_planned 0.0072 (0.0048, 0.0093), brier_skill 0.0054 (0.0033, 0.0079); $RUN6/ablations/backtest-upworthy-20261106-2dc461f52ea4--base-recency-cell-quarter.json brier_skill_planned 0.0076 (0.0047, 0.0107), brier_skill 0.0074 (0.0053, 0.0099) (tables/fmt-base-recency.tex; the two recency rules that were run)
\item Every visible row goes into the pool behind the base rate, which carries no level for a period, whereas the prior works from the level of the most recent quarter that is complete. Recency by itself can therefore account for part of the prior's Brier skill. With a quarter term in the base rate, the prior's skill at planned precision falls from 0.0127 to 0.0072, and its interval stays above zero (Table~\ref{tab:fmt-base-recency}). Only two recency rules for the base rate were tried, so a better rule may take away more.
% src: src/mootcourt/predictors (the boosted model and the local comparison set issue a normal predictive distribution); $RUN6/backtest-upworthy-20261106-2dc461f52ea4.json tables.aggregate[cell=ALL] logscore_vs_p1 rung=GBM -0.0994 (-0.1128, -0.0872), rung=local_eb -0.1276 (-0.1588, -0.1028), rung=text_lr logscore null (the classifier issues a probability and no density)
\item The challengers keep a normal predictive distribution while the rungs of the ladder are Student-t. The log scores of the boosted model and of the local comparison set are therefore not a fair contest with the ladder, and nobody has tried a Student-t forecast on either of them (Section~\ref{sec:method}). The text classifier issues a probability alone, so it has no log score.
% src: $RUN6/horizon-28/by-origin-upworthy-20261114-3fc2dcaeb6f5.json pooled difference 0.1191 (0.0964, 0.1425), leave_one_origin_out[left_out=2013-04-01] difference 0.0029 (-0.0012, 0.0066); $RUN6/horizon-14/backtest-upworthy-20261114-65fd0788a1b7.json tables.aggregate[cell=ALL,rung=P1] brier_skill 0.0065 (0.0035, 0.0097); $RUN6/horizon-28/backtest-upworthy-20261114-3fc2dcaeb6f5.json tables.aggregate[cell=ALL,rung=P1] brier_skill 0.0071 (0.0044, 0.0102) (tables/fmt-horizon.tex); $RUN6/README.md (at the first origin the image-only cell has two visible rows under a 28-day end; the main ledger assumes seven days)
\item The end of each test is assumed, in the main run at seven days after its last arm was created. The reading of Brier skill does not move when the end is assumed at 14 or at 28 days. The lead on the log score does: at 28 days it is 0.119 nats, a figure dominated by one cell of the first origin, where the base rate is fitted on two visible rows. With that origin left out, the interval for the lead spans zero (Section~\ref{sec:ablations}).
\item Two display fields from after launch sit in the retailer archive at the top level of its features, namely the checkpoint count and the final offset measured in days. No model reads them, which the feature allow-list together with its test enforces, and relocating them would force a rewrite of each snapshot identifier that the run folders point to.
% src: src/mootcourt/decision_value.py (a test's value is its best positive observed lift); Section 6.5 ("We did not test this reading")
\item The target of the decision value is the largest observed lift among the arms of a test, and that target rewards noise: when several arms are noisy, the biggest of them comes out bigger through chance alone. The matched baseline strips out the share of that inflation which arm count produces, and we did not test whether what is left inflates a ranking outcome (Section~\ref{sec:decision-results}).
\end{itemize}

% src: src/mootcourt/harness_stats.py _TestSums.resample_counts (one multinomial draw over the scored tests per resample, on scores computed once from the fitted models); Section 5 (intervals resample whole tests)
\paragraph{What the intervals hold fixed.} Every bootstrap interval resamples whole tests as independent units and holds the fitted models and the seven origins fixed. It says how a score would vary over other tests from these same quarters, and as a statement about a new quarter it is likely too narrow. Only the prior's log score lead over the base rate has an interval with the quarters as units, and that interval contains zero (Table~\ref{tab:fmt-by-origin}).

%% tmlr:move-begin morelimits
%% tmlr:note morelimits Limits of the gates, of the interval families and of the sizing evaluation, and the work that remains undone, are in Appendix~\ref{sec:more-limits}.
%% tmlr:title morelimits Further limits
\label{sec:more-limits}
\paragraph{The gates and the families.}
% src: $RUN6/backtest-upworthy-20261106-2dc461f52ea4.json tables.aggregate family_size 52, tables.score_by_year family_size 156, tables.decision_value family_size 6
The gates are a convention for sizing (Section~\ref{sec:protocol}), and each claim stands on its interval. For the headline archive, the aggregate table and the year table are read over 52 and 156 intervals respectively, and adjusted bounds sit next to each. The decision table reads six intervals per budget, while the gate table, the paired ranking table and the sizing table make their adjustments inside the smaller families named in their captions. No adjusted bounds appear in the tables that remain, the calibration tables included, and those are read descriptively.
% src: src/mootcourt/card.py (the released gates are sized for skills of 0.20 and 0.10); $RUN6/paper-extras-upworthy-20261106-2dc461f52ea4.json gates_by_stratum.strata.ALL rungs.P1.verdict "decidable at skill 0.20 and 0.10", rungs.P1.brier_skill 0.0069, rungs.P2.brier_skill 0.0081 (both under 0.01)
The released gates are sized for skills of 0.10 and 0.20. The skills measured here are under 0.01, far below both, so a decidable verdict certifies no skill of the measured size.
% src: $RUN6/backtest-upworthy-20261106-2dc461f52ea4.json header.bootstrap_resamples 500 (a rank among 500 draws moves 0.2 points per draw); $RUN6/paper-extras-upworthy-20261106-2dc461f52ea4.json issued_ece_by_origin null_draws 500, all_rows.by_origin ece_percentile 77.9 and 34.1 at the first two origins (half steps, from the mid-rank rule of Section 11)
The calibration percentile is a rank within 500 draws, so its movement comes in steps of 0.2 points, dropping to half a step whenever a draw matches the observed value, and that explains why certain percentiles as printed finish on an odd tenth (Section~\ref{sec:repro}).

\paragraph{The sizing evaluation.}
% src: docs/superpowers/plans/2026-11-06-x10-sizing-design.md; $RUN6/paper-extras-upworthy-20261106-2dc461f52ea4.json sizing_evaluation.test_and_roll (primary_multiple 20, world_scales; tables/fmt-sizing.tex)
Every test that concluded was run a single time at a single size, so there is no way to replay sizing. Regret is what the evaluation simulates, over a synthetic horizon measured in impressions that is a multiple of the test scale of the cell, with no money involved, inside a world where true effects are fitted using the rows scored at that origin and cell. The outcome appears alongside the pair of alternative worlds in Table~\ref{tab:fmt-sizing}, and for a difference as small as this one another world could flip the sign.

\paragraph{What remains undone.}
\begin{itemize}
\item In the method, seven things:
\begin{itemize}
\item As the realized click-through level goes up, coverage comes down (Section~\ref{sec:calibration-results}). Mixing over the uncertain level inside a forecast might be enough to close the gap.
\item The width of the precedent rung's forecast is not chosen on held-out rows.
\item In the positive control the precedent rung recovers a small share of the planted signal (Appendix~\ref{app:synthetic}). Whether the rest is lost to the guard, to the selection rule, to the kernel or to the brief histories of the early origins is not measured.
\item For the real archive, no held-out split is used to measure what the guard gains inside small cells. That measurement exists for the synthetic ledgers alone (Table~\ref{tab:guard-gain-held-out}).
\item A fixed rule, rather than any estimate from the data, determines how large P1's deviation for a period is.
\item Choosing between the base rate and the prior is a decision no guard makes. The drift rule of the loop stands in for such a guard.
% src: src/mootcourt/config.py NU_GRID (3, 5, 10, 30, inf)
\item The grid of degrees of freedom for the Student-t prior ends at 3. Fits sitting on that edge might want a heavier tail, and we report them just as they stand.
\end{itemize}
\item In the protocol: planned-exposure and planned-duration columns in the exchange format, and adjusted bounds for the tables that lack them.
\item Not run: a run of the challengers on a single headline, so nobody has measured what the reference headline bought the text classifier.
\end{itemize}
%% tmlr:move-end morelimits

\paragraph{Ethics, data norms and misuse.} Each archive is aggregate, comes out under CC BY 4.0 and carries a credit by name. The system takes in no user-level data at all, and we use nothing beyond the exploratory split of the headline archive (Section~\ref{sec:setting}). Launching a change, splitting traffic and stopping a test all lie outside what the system does, but when a team runs just the top of its ranking, the forecast ends up deciding which ideas reach a live test. On this archive the evidence behind that ranking is thin, so dropping a candidate purely on a card's probability is the first misuse we can foresee. Three more points guard against misuse.
\begin{itemize}
\item When a team ranks by win probability, the ledger it writes is selected on what the forecaster itself produced, so each ledger row ought to say whether a card influenced the test, and skill for those rows ought to be reported separately.
\item A team's own forecast must be entered before its card is seen.
\item For a card that abstained, the base rate has to carry a label as the quantity the card declined to improve on.
\end{itemize}
